% =============================================================================
% Lecture 1: Introduction
% Agentic Machine Learning Systems (Volume III), Chapter 1
% =============================================================================
% Thesis: a delegated task is accepted only when a runtime closes the loop
% around the model, bounding what it may do below and checking what it did
% above, in proportion to the task's horizon, state, and authority.
% Running case: the repair agent (Lighthouse 1.1), shown as a real system: one
% recorded repair of astropy issue 14978 from a public SWE-bench Verified run,
% and the 500-task run it belongs to.
%
% Build: make deck V=3 D=01_introduction   (figures: images/make_figs.py,
%        numbers: images/calcs.py, data: images/extract_run.py)
% Plan:  lecture_map.yml in this folder.
% =============================================================================

\documentclass[aspectratio=169, 12pt]{beamer}
\usepackage{../../assets/beamerthememlsys}
\mlsyssetup{
  volume={Volume III},
  chapter={Chapter 1},
  chaptertitle={Introduction},
  layer=trajectory,
  logo={../../assets/img/logo-mlsysbook.png},
  instlogo={../../assets/img/logo-harvard.png},
}
\usepackage[T1]{fontenc}
\usepackage[scaled=0.9]{helvet}
\usepackage{courier}
\renewcommand{\familydefault}{\sfdefault}
\usepackage{booktabs,amsmath}
\graphicspath{{images/}}
% Generated by calcs.py. Do not edit; edit calcs.py and rerun it.
\defval{hook_calls}{28}
\defval{hook_cost}{0.34}
\defval{hook_f2p}{1}
\defval{hook_p2p}{179}
\defval{run_tasks}{500}
\defval{run_submitted}{500}
\defval{run_acc}{324}
\defval{run_rej}{176}
\defval{run_acc_pct}{64.8}
\defval{run_acc_frac}{0.648}
\defval{run_rej_pct}{35}
\defval{run_total}{185.73}
\defval{run_mean}{0.371}
\defval{run_per_acc}{0.57}
\defval{run_per_acc3}{0.573}
\defval{run_winners}{0.320}
\defval{run_understate}{1.8}
\defval{run_understate2}{1.79}
\defval{run_calls_min}{9}
\defval{run_calls_med}{33}
\defval{run_calls_max}{123}
\defval{run_cmds}{18{,}582}
\defval{run_write_pct}{28}
\defval{run_run_pct}{24}
\defval{run_net_word}{nine}
\defval{run_trunc}{298}
\defval{harness_steps}{30}
\defval{harness_window}{128{,}000}
\defval{harness_window_k}{128}
\defval{step_ok}{0.99}
\defval{step_ok_pct}{99}
\defval{step_err_pct}{1}
\defval{h_short}{10}
\defval{h_long}{100}
\defval{h_ext}{200}
\defval{p_short_pct}{90.4}
\defval{p_long_pct}{36.6}
\defval{p_long}{0.366}
\defval{p_ext_pct}{13.4}
\defval{target}{0.9}
\defval{target_pct}{90}
\defval{eps_req}{0.00105}
\defval{eps_req_pct}{0.11}
\defval{eps_ratio}{9.5}
\defval{eps_h}{0.105}
\defval{c_tokens}{1.30}
\defval{c_other}{0.80}
\defval{p_acc}{0.40}
\defval{p_acc_pct}{40}
\defval{p_better}{0.60}
\defval{p_better_pct}{60}
\defval{token_cut}{0.5}
\defval{c_attempt}{2.10}
\defval{attempts}{2.5}
\defval{c_eff}{5.25}
\defval{c_cheap}{1.45}
\defval{c_eff_cheap}{3.63}
\defval{c_eff_better}{3.50}
\defval{token_share}{0.62}
\defval{token_share_pct}{62}
\defval{p_breakeven}{0.58}
\defval{c_understate}{2.5}
\defval{ctx_turns}{30}
\defval{ctx_s0}{8{,}000}
\defval{ctx_delta}{1{,}500}
\defval{ctx_last}{51{,}500}
\defval{ctx_window_pct}{40}
\defval{ctx_total}{892{,}500}
\defval{ctx_ratio}{17.3}
\defval{metr_first_s}{3}
\defval{metr_last_h}{17}
\defval{metr_growth}{19{,}000}
\defval{metr_all_months}{6.2}
\defval{metr_recent_months}{4.2}

\setcounter{mlsystotalsections}{5}

\title{Introduction}
\author{Vijay Janapa Reddi}
\institute{Harvard University}
\date{}

% A large answer choice for predict and concept-question slides
\newcommand{\choice}[2]{\textbf{\textcolor{mlsysaccent}{#1}}\enspace #2}
% A choice highlighted from overlay #1 on (the answer)
\newcommand{\answerat}[3]{\alt<#1->{\colorbox{mlsysaccent!15}{\choice{#2}{#3}}}{\choice{#2}{#3}}}
% A goal or takeaway line with its thumbnail
\newcommand{\thumbline}[2]{%
  \begin{minipage}[c]{2.6cm}\centering\includegraphics[width=2.5cm,height=0.95cm,keepaspectratio]{#1}\end{minipage}%
  \hspace{0.3cm}\begin{minipage}[c]{0.74\textwidth}\raggedright\small #2\end{minipage}}
% One line of a worked example: shown from overlay #1 on
\newcommand{\wstep}[2]{\uncover<#1->{#2}\\[0.2cm]}
% Inline code
\newcommand{\code}[1]{{\ttfamily\small #1}}
% Golden-thread tag: one line on the goals slide naming the H·S·A term this lecture acts on
\newcommand{\hsathread}[1]{{\scriptsize\textbf{\textcolor{mlsysaccent}{H·S·A thread:}} #1\par}}

\begin{document}

% =============================================================================
% TITLE AND HOOK
% =============================================================================
\mlsysnote{[0.5 min] [CORE] Read the thesis aloud once. This lecture defines the volume's tools: the closed loop, the trajectory, the H·S·A exposures, invariant closure, the task contract, and the two relations every later lecture prices with. One real system carries it: a recorded repair agent fixing a real bug, and the \val{run_tasks}-task run it came from. The notes total about 71 minutes, of which about 61 are [CORE]. FLEX: cutting every [OPTIONAL] frame leaves a complete lecture; cut first the lighthouse cards, the workflow-versus-loop frame, and the open questions.}
\mlsystitle{Introduction}{A delegated task is accepted only when a runtime closes the loop around the model, bounding what it may do and checking what it did.}{cover_introduction_blueprint.png}

\begin{frame}{Who decided this repair was accepted?}
  \mlsysactivity{Vote}{2 min}
  \begin{columns}[c]
    \begin{column}{0.56\textwidth}
      \centering\safeimg[width=\textwidth,height=4.6cm,keepaspectratio]{issue_14978.png}
    \end{column}
    \begin{column}{0.42\textwidth}
      \small A repair agent took this real issue, worked \textbf{\val{hook_calls} model calls}, spent \textbf{\$\val{hook_cost}}, and ended: \textit{``The fix is working perfectly.''}\\[0.25cm]
      \choice{A}{its final message}\\[0.08cm]
      \choice{B}{its own test scripts}\\[0.08cm]
      \answerat{2}{C}{tests it never saw}\\[0.08cm]
      \choice{D}{a human reviewer}
    \end{column}
  \end{columns}
  \mlsyscredit{GitHub, astropy/astropy issue 14978 (screenshot). Run: SWE-bench Verified, astropy\_\_astropy-14995, mini-swe-agent 1.0.0 with Claude 4 Sonnet (SWE-bench experiments, 2025).}
  \note{[2.5 min] [CORE] NARRATE: a real bug report in astropy, a scientific Python library: multiplying data with a mask by data without one crashes. A repair agent, a model in a loop that runs shell commands in a copy of the repository, read the issue, made \val{hook_calls} model calls, spent \$\val{hook_cost}, and ended with the sentence on the slide. ENGAGE: vote by hand. Step 2: C. The benchmark applied the agent's patch and ran a test the agent never saw, \code{test\_nddata\_bitmask\_arithmetic}, plus \val{hook_p2p} existing tests. Only then did the task count as resolved. Students who pick A trust the report; B is evidence the agent can shape. WARN: the agent's final message was right this time; in the same run \val{run_rej} of \val{run_tasks} equally confident submissions were wrong (segment 3). LINK: the whole lecture explains why a correct-looking action needs a runtime and outside evidence. FLEX: [CORE].}
\end{frame}

\begin{frame}{Which cleanup run is dangerous?}
  \mlsysactivity{Predict}{2 min}
  \centering
  \mlsyssteps[width=0.97\textwidth,height=4.8cm,keepaspectratio]{cleanup_step}{2}\\[0.12cm]
  {\normalsize\choice{A}{run A}\qquad\answerat{2}{B}{run B}\qquad\choice{C}{both}\qquad\choice{D}{neither: the model decides}}
  \mlsyscredit{Introduction, Example 1.1 (illustrative command).}
  \note{[2.5 min] [CORE] NARRATE: the chapter's Example 1.1. The same model, prompt, and horizon issue the same cleanup command that deletes files. In run A the files sit in a disposable clone of the repository; in run B they sit in the developer's only workspace, which holds edits never committed. ENGAGE: vote, then ask whether the model could tell the two runs apart from its prompt. Step 2: B. In A the copy is task-local working state ($S_2$) and the worst damage is a discarded sandbox ($A_1$). In B the same command reaches durable shared state ($S_3$) and destroys work with no snapshot ($A_3$). WARN: students answer D, that a careful model would avoid it. The model sees the same text in both runs; only the environment differs, so only the runtime can tell them apart. Version control does not help: uncommitted edits have no snapshot. LINK: segment 2 names these levels. FLEX: [CORE].}
\end{frame}

\begin{frame}{Recall: what fails on a \val{h_long}-step task?}
  \mlsysactivity{Retrieval}{1 min}
  \centering
  {\Large\bfseries Question of the day}\\[0.4cm]
  {\normalsize Each step is right \val{step_ok_pct}\% of the time, and nothing checks the steps.\\ How often does a \val{h_long}-step task succeed? What is the course's unit of engineering?}\\[0.5cm]
  \uncover<2->{\normalsize\textbf{\textcolor{mlsysaccent}{\val{p_long_pct}\%}}, and the unit of engineering is the \textbf{\textcolor{mlsysaccent}{trajectory}} of one whole task.}
  \note{[1 min] [CORE] LINK: Lecture 0 asked students to predict this; today derives it. Give 40 seconds. Step 2: \val{step_ok} to the \val{h_long}th power is \val{p_long}. The unit of engineering is the trajectory: the whole run of one delegated task, from its first context to acceptance or stop. ENGAGE: ask who remembers what checks do to that number; segment 4 answers. FLEX: [CORE]; a class that skipped Lecture 0 needs only this slide.}
\end{frame}

\begin{frame}{Why can a correct action still fail the task?}
  \thumbline{thumb_loop_step2.pdf}{Explain why agency is a property of the closed loop around the model, and trace one trajectory.}\\[0.2cm]
  \thumbline{thumb_hsa_step3.pdf}{Classify a task by horizon, state, and authority, and name what each exposure demands.}\\[0.2cm]
  \thumbline{thumb_ladder_step2.pdf}{Separate bounds enforced below the model from evidence checked above it; write a task contract.}\\[0.2cm]
  \thumbline{thumb_horizon_step3.pdf}{Compute unchecked success over a horizon and cost per accepted task, and trace a turn's six phases.}\\[0.25cm]
  \hsathread{defines all three exposures and invariant closure}
  \note{[1 min] [CORE] The driving question, from the chapter's Purpose: why can a model that generates a fluent, correct action still fail to complete an autonomous task? Four goals, one per thumbnail. Segment 1 is the loop, segment 2 the exposures, segment 3 closure and the contract, segments 4 and 5 the arithmetic and the turn. Point out that every segment returns to the astropy repair.}
\end{frame}

% =============================================================================
% SECTION: CLOSED LOOP
% =============================================================================
\mlsysnote{[0.5 min] [CORE] First: who closes the loop, and what authority does the model hold?}
\mlsysdivider{Who closes the loop, and what does the model hold?}{Introduction, \S1.1 to \S1.3}
\section{Closed Loop}

\begin{frame}[fragile]{A real agent harness queries, runs, appends, and repeats}
  \begin{columns}[c]
    \begin{column}{0.60\textwidth}
\begin{semiverbatim}\footnotesize
while True:
    self.step()          \textcolor{gray}{# one turn}
    if self.messages[-1]
         .get("role") == "exit":
        break

def step(self):
    return self.execute_actions(
             self.\textcolor{BlueLine}{\textbf{query}}())

outputs = [self.env.\textcolor{OrangeLine}{\textbf{execute}}(action)
           for action in actions]
\end{semiverbatim}
    \end{column}
    \begin{column}{0.37\textwidth}
      \small\textcolor{BlueLine}{\textbf{query}}: the model returns text with one proposed command.\\[0.3cm]
      \textcolor{OrangeLine}{\textbf{execute}}: software runs it and appends the output.\\[0.3cm]
      No person sits inside the loop.
    \end{column}
  \end{columns}
  \mlsyscredit{Excerpt, mini-swe-agent 2.4.6, the current version; the 2025 runs in this lecture used 1.0.0. agents/default.py (commit 04d809c, MIT); exception handling omitted.}
  \note{[2 min] [CORE] NARRATE: real code from a small open-source repair agent of the kind that fixed the astropy bug. This is its current version, 2.4.6; the 2025 runs used version 1.0.0, whose loop has the same shape and stops by raising an exception instead of appending an exit message. The loop calls \code{step}; \code{step} queries the model and executes whatever actions come back; the outputs become messages the next query reads. The loop ends only when a message with the role \code{exit} appears. ENGAGE: where in these lines does anything check whether the command is allowed, or whether the task is done? (Expected: nowhere here; the environment object and the exit rule hold those decisions, and segments 2 and 5 open them.) WARN: students read ``agent'' as a property of the model. The model call is one line; everything else is ordinary software. FLEX: [CORE].}
\end{frame}

\begin{frame}{In an agent, a runtime closes the loop a person used to close}
  \centering
  \mlsyssteps[width=0.97\textwidth,height=5.8cm,keepaspectratio]{loop_step}{2}
  \mlsyscredit{Redrawn from Introduction, figure 1.}
  \note{[2 min] [CORE] Step 1: a conversational assistant. The model answers; a developer runs every command, reads every error, and pastes it back. The person closes the loop and carries the memory. (A predictive service, the earlier arrangement, has no loop at all: application code consumes one score.) Step 2: an agentic system replaces the person with a runtime. The runtime assembles context, receives a proposal, dispatches only actions it grants into a sandbox, and takes its evidence from verifiers, treating the model's account as a claim. The dashed line separates proposing from acting. ENGAGE: in panel B, which box can delete a file? (Expected: only the sandbox, and only when the runtime dispatches.) WARN: agency is a property of the closed loop. A more capable model proposes better actions; permissions, isolation, persistence, and completion checks stay in the runtime. FLEX: [CORE].}
\end{frame}

\begin{frame}{The stochastic computer maps an agent onto a computer}
  \begin{columns}[c]
    \begin{column}{0.64\textwidth}
      \centering\mlsyssteps[width=\textwidth,height=5.6cm,keepaspectratio]{stochastic_step}{2}\\
      \uncover<2->{\footnotesize\textcolor{RedLine}{\textbf{Unlike a processor, the model only proposes.}}}
    \end{column}
    \begin{column}{0.34\textwidth}
      \centering
      \safeimg[width=\textwidth,height=2.9cm,keepaspectratio]{edvac.jpg}\\[0.05cm]
      {\footnotesize EDVAC: its control unit executes every instruction it fetches.}\\[0.1cm]
      \safeimg[width=0.48\textwidth,height=1.7cm,keepaspectratio]{edvac_first_draft.jpg}
    \end{column}
  \end{columns}
  \mlsyscredit{Map redrawn from Introduction, figure 2. Photos: U.S. Army, EDVAC at BRL (public domain); von Neumann, First Draft of a Report on the EDVAC, 1945 (public domain); Wikimedia Commons.}
  \note{[3 min] [CORE] Shown once in this course. Step 1: an agentic system needs the parts of a computer. The model plays the processor, context and stores the memory hierarchy, tools input and output, the runtime the operating system. Across tasks, learning from recorded runs plays recompiling, and many agents play scale-out. The photo is EDVAC, whose 1945 design report by von Neumann split a computer into arithmetic, control, memory, and input-output; its control unit held the authority to execute. Step 2: the map breaks at the processor. A processor executes deterministically; a model samples a proposal that can be fluent and wrong, and cannot vouch for it. So authority inverts: the runtime decides what runs. Memory breaks too: context is a lossy copy of the real state in files and databases. ENGAGE: which classical guarantee survives? (Expected: access control and logging, enforced by the runtime.) WARN: after today we name parts by their own names, never ``processor'' for the model. FLEX: [CORE].}
\end{frame}

\begin{frame}{A trajectory records proposals, decisions, and observations}
  \centering
  {\normalsize$\tau = (s_0, a_0, o_1, s_1, a_1, o_2, \dots, s_H)$}\quad{\footnotesize state $s$, proposed action $a$, observation $o$; $H$ turns}\\[0.12cm]
  \mlsyssteps[width=0.95\textwidth,height=5.2cm,keepaspectratio]{traj_step}{3}
  \mlsyscredit{Verbatim excerpts, turns 11 to 16 of the astropy\_\_astropy-14995 trajectory (SWE-bench experiments, mini-swe-agent 1.0.0, Claude 4 Sonnet). Notation: Introduction, Definition 1.2.}
  \note{[2.5 min] [CORE] Definition 1.2: a trajectory is the ordered record of states, proposed actions, runtime decisions, and observations of one delegated task, from its first context to acceptance, termination, or escalation. One turn is one model call plus the runtime's handling of its proposal. Step 1: the real run. Turn 11 reproduces the bug: the script reports the crash. Turn 13 proposes a one-line edit; the runtime runs it in the container. Step 2: turn 14 shows the edit broke the file, an indentation error. Turn 16 restores the file with \code{git checkout}. Step 3: the record keeps the failed patch. ENGAGE: if we saved only the final diff, what would we lose? (Expected: the cost of turns 13 to 16, the error, and the evidence for acceptance.) WARN: students equate the transcript of tokens with the trajectory; the trajectory also records what the runtime allowed and what the tools returned. FLEX: [CORE].}
\end{frame}

\begin{frame}{\val{run_tasks} real repairs ran from \val{run_calls_min} to \val{run_calls_max} model calls each}
  \centering
  \safeimg[width=0.88\textwidth,height=4.6cm,keepaspectratio]{calls_hist.pdf}\\[0.05cm]
  {\small Cost, context, and permissions accrue over the whole trajectory.}
  \mlsyscredit{SWE-bench Verified, \val{run_tasks} tasks, one attempt each: mini-swe-agent 1.0.0 with Claude 4 Sonnet (SWE-bench experiments, per\_instance\_details.json). Counts: calcs.py.}
  \note{[1.5 min] [CORE] NARRATE: each bar counts tasks by how many model calls their trajectory took, green accepted and red rejected by the hidden tests. The median task took \val{run_calls_med} calls; the longest took \val{run_calls_max}. ENGAGE: what does a serving system that bills and forgets per request miss here? (Expected: the task's running total of cost, the state carried between calls, and the permissions that must hold for every call.) WARN: the red share grows to the right, but harder tasks also take more calls, so this plot does not by itself prove that length causes failure; segment 4 does the arithmetic. FLEX: [CORE]; the segment's real data.}
\end{frame}

\begin{frame}{A stronger model: which job moves into it?}
  \mlsysactivity{Concept question}{1.5 min}
  \centering
  {\normalsize A team replaces its model with one that proposes better patches.\\ Which runtime job can it now drop?}\\[0.5cm]
  {\normalsize\choice{A}{permission checks}\qquad\choice{B}{crash recovery}\qquad\choice{C}{completion checks}\qquad\answerat{2}{D}{none}}\\[0.4cm]
  \uncover<2->{\small A better model raises acceptance. Enforcement, isolation, persistence, and checks stay in the runtime.}
  \note{[1.5 min] [CORE] ENGAGE: vote, pair for 30 seconds, vote again. Step 2: D. A stronger model makes fewer wrong proposals, which raises the share of attempts accepted; it still holds no authority, keeps no durable state, and cannot certify its own result. Students who pick C argue the better model is more often right; ask how the runtime would know which runs are right without a check. Land the segment aloud: agency is a property of the closed loop, and we account for cost and success over the whole trajectory. FLEX: [CORE].}
\end{frame}

% =============================================================================
% SECTION: EXPOSURE
% =============================================================================
\mlsysnote{[0.5 min] [CORE] Second: how far does a task depart from a single model call, and what does each departure demand?}
\mlsysdivider{How far does a task depart from one model call?}{Introduction, \S1.4}
\section{Exposure}

\begin{frame}{Three axes measure how far a task departs from one call}
  \centering
  \mlsyssteps[width=0.97\textwidth,height=5.9cm,keepaspectratio]{hsa_step}{3}
  \mlsyscredit{Redrawn from Introduction, figure 3; positions are illustrative task configurations.}
  \note{[2.5 min] [CORE] Definition 1.3, the H·S·A exposure space. Step 1: the origin is one model call: no turns, no carried state, no effects. Horizon $H$ counts turns: each one is another chance to act on a stale or wrong observation. Step 2: state runs from $S_0$ stateless, $S_1$ the context and serving state, $S_2$ a task-local working copy, to $S_3$ durable shared state that outlives the task. Authority runs from $A_0$ read-only to $A_3$ irreversible. Step 3: the four lighthouse workloads. The repair agent: about 100 turns, a working copy, workspace edits. The remediation agent: about 20 turns but $S_3$ and $A_3$. ENGAGE: which agent needs the strongest approval gate, and which the strongest checks over a long run? (Expected: remediation; repair or research.) WARN: the coordinate is where a task sits; closure is what the runtime builds in answer. FLEX: [CORE].}
\end{frame}

\begin{frame}{The consequence in its environment sets an action's authority}
  \centering
  \safeimg[width=0.92\textwidth,height=5.4cm,keepaspectratio]{authority.pdf}
  \mlsyscredit{Introduction, \S1.4 and Example 1.1; Lighthouse 1.4 for the seat hold and refund.}
  \note{[2 min] [CORE] NARRATE the four side-effect classes bottom up. $A_0$ reads only, yet an agent that reads a secret can leak it in its output. $A_1$ mutates a disposable workspace: tearing it down undoes everything. $A_2$ reaches external systems, but every effect pairs with an executable compensating action, run under checked preconditions; releasing a seat hold is the support agent's example. $A_3$ has no automated reversal: a refund paid out, an unversioned production record. Then the bottom line: the hook's delete is $A_1$ in the clone and $A_3$ on the only copy. ENGAGE: is \code{git push} to a shared branch $A_1$, $A_2$, or $A_3$? (Expected: it depends on whether a reviewed revert exists and on who pulls it; argue from consequence.) WARN: students classify by verb. LINK: Lecture 12 teaches compensation; Lecture 7 sets the dispatch policy for each class. FLEX: [CORE].}
\end{frame}

\begin{frame}[fragile]{The environment class picked at launch grants the authority}
  \begin{columns}[T]
    \begin{column}{0.49\textwidth}
      {\small\textbf{Local: runs on the host}}
\begin{semiverbatim}\fontsize{9}{10.6}\selectfont\bfseries
process = subprocess.Popen(
    command,
    \textcolor{RedLine}{\textbf{shell=True}},
    cwd=cwd,
    \textcolor{RedLine}{\textbf{env=os.environ | ...}},
    ...)
\end{semiverbatim}
      {\footnotesize Host files, host credentials.}
    \end{column}
    \begin{column}{0.49\textwidth}
      {\small\textbf{Docker: runs in a container}}
\begin{semiverbatim}\fontsize{9}{10.6}\selectfont\bfseries
cmd = [\textcolor{GreenLine}{\textbf{"docker", "exec"}}, "-w", cwd]
cmd.extend([self.container_id,
            *self.config.interpreter,
            command])
...
\textcolor{GreenLine}{\textbf{docker rm -f}} \{container_id\}
\end{semiverbatim}
      {\footnotesize A copy that cleanup throws away.}
    \end{column}
  \end{columns}
  \vspace{0.25cm}
  \centering{\small With the same model and loop, the class picked at launch sets the authority level.}
  \mlsyscredit{Excerpts, mini-swe-agent 2.4.6, the current version (the 2025 runs used 1.0.0): environments/local.py and docker.py (commit 04d809c, MIT).}
  \note{[2 min] [CORE] NARRATE: the same agent ships two environment classes. The local one hands the model's command to a shell on the host, with the host's working directory and every environment variable, including any API keys. The Docker one runs the command inside a container and deletes the container at cleanup. The prompt can be identical. This is Example 1.1 in real code: the environment decides between $A_1$ and something close to $A_3$. Zero ambient authority means the model holds no standing privilege; every effect passes an explicit grant the runtime checks. ENGAGE: in the local class, what stops the model from reading \code{\textasciitilde/.ssh}? (Expected: nothing in this code.) WARN: a container is a boundary only if its mounts and network are restricted; Lecture 8 builds containment beneath the model. FLEX: [CORE].}
\end{frame}

\begin{frame}{\val{run_cmds} real repair commands read, wrote, or ran code in a copy}
  \centering
  \safeimg[width=0.90\textwidth,height=3.7cm,keepaspectratio]{cmd_classes.pdf}\\[0.1cm]
  {\small Every write landed in a disposable container, so authority stayed at $A_1$.}
  \mlsyscredit{Same SWE-bench Verified run, all \val{run_tasks} trajectories; command classes by pattern (extract\_run.py, derived).}
  \note{[1 min] [OPTIONAL] NARRATE: every command the agent proposed across the \val{run_tasks} tasks, sorted by what it does. Almost half only read files; \val{run_write_pct} percent changed files in the working copy; \val{run_run_pct} percent executed code there. Only \val{run_net_word} tried the network, for instance a package install. All of it happened inside throwaway containers, which is why the run could let the model act freely: the worst outcome was a discarded copy. ENGAGE: which of these classes would change level if the same agent ran with the local environment? (Expected: writes and runs, which would reach the host.) WARN: the classes are assigned by matching command patterns, so a few are misfiled; the shares are approximate. FLEX: cut if short; the previous slide carries the point.}
\end{frame}

\begin{frame}{Four lighthouse workloads bind on four constraints}
  \begin{columns}[c]
    \begin{column}{0.64\textwidth}
      \footnotesize
      \begin{tabular}{@{}P{1.9cm}P{2.6cm}P{4.1cm}@{}}\toprule
        \textbf{Lighthouse} & \textbf{H·S·A} & \textbf{Binding constraint}\\\midrule
        \textcolor{mlsysaccent}{\textbf{Repair}} & 100, $S_2$, $A_1$ & errors compound; history resent each turn\\[0.04cm]
        Research & 1,000, $S_3$, $A_0$ & untrusted pages can steer it\\[0.04cm]
        Remediation & 20, $S_3$, $A_3$ & one wrong change costs an outage\\[0.04cm]
        Support & 3, $S_1$+$S_3$, $A_0$--$A_3$ & tail latency for waiting users\\\bottomrule
      \end{tabular}
    \end{column}
    \begin{column}{0.34\textwidth}
      \centering\mlsysphoto[width=\textwidth,height=3.6cm,keepaspectratio]{datacenter-server-racks.jpg}\\[0.05cm]
      {\footnotesize Shared infrastructure: where the remediation agent acts.}
    \end{column}
  \end{columns}
  \mlsyscredit{Introduction, \S1.4.1, Lighthouses 1.1 to 1.4. Photo: Carl Lender, CC BY 2.0, Wikimedia Commons.}
  \note{[1.5 min] [OPTIONAL] The volume's four recurring workloads. The repair agent is the running example in every lecture: \val{harness_steps} steps budgeted in a \val{harness_window}-token window, bound by compounding errors and by input that grows when every turn resends the history. The research agent only reads, yet a retrieved page can carry text the model follows as an instruction, prompt injection, which Lecture 9 owns. The remediation agent changes live production systems; approval latency and blast radius dominate. The support agent serves thousands of short sessions with a user waiting. ENGAGE: why does a mechanism that helps one lighthouse hurt another? (Expected: summarizing context helps a 100-turn repair and only adds latency to a 3-turn support session.) FLEX: cut if short.}
\end{frame}

\begin{frame}{Which exposure moved, and what closes it?}
  \mlsysactivity{Diagnose}{2.5 min}
  \small
  The repair agent edits a disposable working copy and returns a patch for review.\\[0.25cm]
  \begin{tabular}{@{}P{6.3cm}P{7.6cm}@{}}\toprule
    \textbf{Change to the task} & \leavevmode\uncover<2->{\textbf{Exposure, and what closes it}}\\\midrule
    1. The task now includes deployment. & \leavevmode\uncover<2->{Authority rises to $A_3$; a human approval no test can supply}\\[0.08cm]
    2. The copy holds the only copy of unrelated edits. & \leavevmode\uncover<2->{State rises to $S_3$; grant an isolated clone instead}\\[0.08cm]
    3. A sealed sequential test passes on the race fix. & \leavevmode\uncover<2->{Evidence covers one thread; add multi-threaded stress runs}\\\bottomrule
  \end{tabular}
  \mlsyscredit{Introduction, Checkpoint 1.1.}
  \note{[2.5 min] [CORE] ENGAGE: pairs take the three rows for 90 seconds. Step 2: reveal. Row 1: deployment is irreversible authority; the contract must add an approval decision, and no passing test substitutes for it. Row 2: the workspace is now durable shared state; the single change is to grant the agent an isolated clone and integrate reviewed patches through a separate gate. Row 3: a sealed sequential test protects its assertions but covers single-threaded runs only; a race fix needs stress runs across thread schedules, then human review. WARN: students fix row 2 with a prompt telling the model to be careful; the host still allows what its credentials allow. Land the segment aloud: horizon, state, and authority locate a task, authority comes from consequence in the environment, and the runtime sizes its defenses to that location. FLEX: [CORE].}
\end{frame}

% =============================================================================
% SECTION: CLOSURE
% =============================================================================
\mlsysnote{[0.5 min] [CORE] Third: the model cannot vouch for its own output. What keeps a plausible wrong answer out?}
\mlsysdivider{How does the runtime keep a plausible wrong answer out?}{Introduction, \S1.5}
\section{Closure}

\begin{frame}{A fail-plausible fault looks like a success}
  \centering
  \mlsyssteps[width=0.95\textwidth,height=4.0cm,keepaspectratio]{faults_step}{3}
  \mlsyscredit{After Introduction, \S1.5.1 margin figure; Schlichting and Schneider 1983; Lamport et al. 1982; Avizienis et al. 2004.}
  \note{[2 min] [CORE] Step 1: a fail-stop component halts, so a supervisor sees a missing heartbeat or a nonzero exit code. Step 2: a Byzantine component may behave arbitrarily, and agreement among $3f + 1$ replicas tolerates $f$ such faults. Step 3: the agent's characteristic failure is neither. A patch compiles yet deletes the lock that made the code thread safe; a test run exits 0 because the failing assertion was commented out; a summary reports success. The model process did not crash, so crash monitors log a success; samples from one model share blind spots, so they can agree on the wrong answer. In dependability terms it is an unsignaled content failure. ENGAGE: name a check that would catch the deleted lock. (Expected: a concurrency stress test the agent cannot edit.) WARN: more samples from the same model is not redundancy. FLEX: [CORE].}
\end{frame}

\begin{frame}[fragile]{All \val{run_submitted} runs claimed success; hidden tests rejected \val{run_rej}}
  \begin{columns}[c]
    \begin{column}{0.53\textwidth}
      {\footnotesize The loop's exit rule (\code{agents/default.py}):}
\begin{semiverbatim}\scriptsize
if lines and lines[0].strip() ==
   "MICRO_SWE_AGENT_FINAL_OUTPUT":
    raise Submitted(...)
\end{semiverbatim}
      {\footnotesize The model prints a word; the loop stops.\\ Acceptance happens elsewhere.}
    \end{column}
    \begin{column}{0.45\textwidth}
      \centering\safeimg[width=\textwidth,height=4.3cm,keepaspectratio]{claimed.pdf}
    \end{column}
  \end{columns}
  \mlsyscredit{Code: mini-swe-agent 1.0.0, the version that made these runs, agents/default.py lines 126 and 127 (commit a2b0efee5f, MIT). Data: SWE-bench Verified run as before; every trajectory ended with exit status Submitted.}
  \note{[2 min] [CORE] NARRATE: the stop rule of the real loop. When the first line of a command's output is a fixed marker, the loop records the task as submitted. That is the model's own report. In the run, all \val{run_submitted} tasks ended that way; the benchmark's hidden tests then accepted \val{run_acc} and rejected \val{run_rej}, \val{run_rej_pct} percent. ENGAGE: if a runtime accepted on the submission alone, what share of its acceptances would be wrong? (Expected: \val{run_rej} of \val{run_tasks}.) WARN: this harness is honest about it: it is built for a benchmark whose grader sits outside the loop. Copying the loop into production without the grader copies only the self-report. LINK: this is the chapter's fallacy that a coherent report establishes completion. FLEX: [CORE]; the segment's real data.}
\end{frame}

\begin{frame}{Five evidence levels support five different claims}
  \centering
  \mlsyssteps[width=0.97\textwidth,height=5.7cm,keepaspectratio]{ladder_step}{2}
  \mlsyscredit{Introduction, \S1.5.2. Examples: the astropy\_\_astropy-14995 trajectory and its SWE-bench Verified record (FAIL\_TO\_PASS and PASS\_TO\_PASS tests).}
  \note{[2.5 min] [CORE] Step 1: the closure evidence levels, weakest by default at the bottom. Self-report is the model's claim about the environment, with no observation behind it. Static analysis checks types and structure without running. Visible tests run cases the agent can read, and so can fit. Sealed tests are withheld and tamper-proof. Formal verification proves the code meets a written specification, which may itself be wrong. Step 2: the astropy run used three. The agent's sentence is level 1. Its run of the existing test file is level 3. The benchmark's new test, never shown to the agent, plus the \val{hook_p2p} old ones, is level 4. ENGAGE: why does the order say ``by default''? (Expected: coverage decides; a thorough sealed suite beats a proof of an incomplete specification.) WARN: every check is a bounded claim narrower than the task: the accepted diff also carried four scratch scripts no test looked at. A grader's own false-acceptance rate is measured in Lecture 13. FLEX: [CORE].}
\end{frame}

\mlsysnote{[1.5 min] [CORE] The volume's one principle defined in a chapter. Read the two halves slowly. Below the model: between its proposal and the environment, where the runtime decides what may execute: a path check, a container, a capability grant. Above the model: between the finished result and whoever accepts it, where evidence decides what counts as done. The formula is the price of neither: with $H$ independent unchecked steps, each wrong with probability $\epsilon$, success is $(1-\epsilon)^H$. A bound rejects the violations it blocks; a check rejects those it detects. ENGAGE: in the astropy run, name one bound and one check. (Expected: the container; the hidden test.) WARN: greater exposure demands stronger bounds and stronger evidence, and the contract itself is adopted by people: the model may draft it but never adopt or relax it.}
\mlsysfocus{Invariant closure}{%
{\large Constraints never rest on the model.}\\[0.35cm]
{\normalsize Bounds close \textbf{below} it; acceptance rests on evidence \textbf{above} it.}\\[0.5cm]
{\normalsize Without either: \quad $\Pr(\text{success}) = (1-\epsilon)^{H}$}%
}

\begin{frame}{A task contract turns a request into rules the runtime enforces}
  \footnotesize
  \renewcommand{\arraystretch}{1.3}
  \begin{tabular}{@{}P{3.5cm}P{5.4cm}P{5.0cm}@{}}\toprule
    \textbf{Part} & \textbf{In the repair task} & \textbf{If it is left out}\\\midrule
    Goal $G$ & fix the race; leave other behavior alone & delete the failing assertion \\[0.05cm]
    \leavevmode\uncover<2->{Environment $\mathcal{E}$} & \leavevmode\uncover<2->{source baseline, isolated working copy} & \leavevmode\uncover<2->{edits pollute the shared host}\\[0.05cm]
    \leavevmode\uncover<2->{Permitted actions $\mathcal{A}_{\text{perm}}$} & \leavevmode\uncover<2->{edit allowed paths, run tests; no network} & \leavevmode\uncover<2->{host credentials allow everything}\\[0.05cm]
    \leavevmode\uncover<3->{Observations $\mathcal{O}_{\text{avail}}$} & \leavevmode\uncover<3->{source excerpts, bounded logs, diffs} & \leavevmode\uncover<3->{a huge log floods the context}\\[0.05cm]
    \leavevmode\uncover<3->{Completion criteria $\mathcal{K}_{\text{comp}}$} & \leavevmode\uncover<3->{sealed suites plus review for the rest} & \leavevmode\uncover<3->{the model's ``done'' ends the run}\\\bottomrule
  \end{tabular}\\[0.12cm]
  \uncover<3->{\small\textcolor{mlsysaccent}{\textbf{A prompt states intent to the model; a contract is enforced by the runtime.}}}
  \mlsyscredit{Introduction, \S1.5.3, Definition 1.4 and Table 2.}
  \note{[2.5 min] [CORE] Definition 1.4: the task contract $\mathcal{C} = \langle G, \mathcal{E}, \mathcal{A}_{\text{perm}}, \mathcal{O}_{\text{avail}}, \mathcal{K}_{\text{comp}} \rangle$. Step 1: the goal carries a negative scope, what must stay untouched; without it, deleting the assertion ``fixes'' the bug. Step 2: the environment and the permitted actions; without them, edits and package installs pollute shared state and tool calls run with whatever the host process holds. Step 3: observations and completion criteria. The real run truncated \val{run_trunc} over-long outputs with a warning, a bounded observation; truncation must be marked, or the model reads a missing error as no error. Without completion criteria the runtime stops on the model's word. Budgets travel with the contract: turns, wall clock, dollars, and what happens when one runs out. ENGAGE: which part decides the denominator of cost per accepted task? (Expected: completion criteria.) FLEX: [CORE].}
\end{frame}

\begin{frame}{True or false?}
  \mlsysactivity{Concept question}{2.5 min}
  \small
  \begin{tabular}{@{}P{6.4cm}P{7.5cm}@{}}\toprule
    \textbf{Claim} & \leavevmode\uncover<2->{\textbf{Verdict}}\\\midrule
    1. A coherent final report shows the task is done. & \leavevmode\uncover<2->{\textbf{False.} \val{run_rej} of \val{run_tasks} confident submissions failed.}\\[0.1cm]
    2. ``Edit only this folder'' in the prompt limits writes. & \leavevmode\uncover<2->{\textbf{False.} The host allows what its credentials allow.}\\[0.1cm]
    3. A \val{harness_window_k}K window removes the need to manage state. & \leavevmode\uncover<2->{\textbf{False.} \val{ctx_last} tokens fit, but stale edits sit beside new ones; \val{ctx_total} tokens are resent.}\\\bottomrule
  \end{tabular}
  \mlsyscredit{Introduction, \S1.9; run data as before.}
  \note{[2.5 min] [CORE] ENGAGE: hands up for true on each, then reveal. 1: false; the real run, and the chapter's arithmetic: an unchecked \val{h_long}-turn repair succeeds \val{p_long_pct} percent of the time, so a runtime that trusts every confident report accepts mostly wrong work. 2: false; an instruction expresses intent, and injected text in a source file can override it. The boundary must live in the execution layer. 3: false; a \val{ctx_turns}-turn trajectory starting at \val{ctx_s0} tokens and adding \val{ctx_delta} per turn ends with a \val{ctx_last}-token call, \val{ctx_window_pct} percent of the window, yet resends \val{ctx_total} tokens in all, \val{ctx_ratio} times the final context, and nothing marks which file version is current. Lecture 2 prices that resend. Land the segment aloud: bounds below the model, evidence above it, both written into a contract before the loop runs. FLEX: [CORE].}
\end{frame}

% =============================================================================
% SECTION: HORIZON AND COST
% =============================================================================
\mlsysnote{[0.5 min] [CORE] Fourth: how do reliability and cost behave over a whole trajectory?}
\mlsysdivider{How do reliability and cost scale with a trajectory?}{Introduction, \S1.6}
\section{Horizon \& Cost}

\begin{frame}{Delegated tasks grew from \val{metr_first_s} seconds to \val{metr_last_h} hours}
  \centering
  \safeimg[width=0.95\textwidth,height=5.3cm,keepaspectratio]{metr.pdf}\\[0.05cm]
  {\small Each doubling of the task needs half the error per step, or checks twice as often.}
  \mlsyscredit{Kwa et al., ``Measuring AI Ability to Complete Long Software Tasks,'' arXiv:2503.14499; METR Time Horizon v1.1 (via MLSysIM). Fit and doubling times: calcs.py, as in Introduction, figure 4.}
  \note{[2 min] [CORE] NARRATE: METR timed experts on software and research tasks, then measured, for each model, the expert-time length of task its agent completes in half its attempts. Each point is one model at its release date; filled points were the frontier; bars are 95 percent intervals. The frontier rose from about \val{metr_first_s} seconds in 2019 to about \val{metr_last_h} hours in 2026, about \val{metr_growth} times, doubling every \val{metr_all_months} months overall and every \val{metr_recent_months} months since 2023. ENGAGE: at fixed accuracy per step, what happens to success when the task doubles? (Expected: it squares; next slide.) WARN: ``50 percent horizon'' means half the attempts at that length still fail; the number measures how long a task can be, and says little about reliability. FLEX: [CORE]; the segment's real data.}
\end{frame}

\begin{frame}{Horizon compounding decays unchecked success exponentially}
  \begin{columns}[c]
    \begin{column}{0.68\textwidth}
      \centering\mlsyssteps[width=\textwidth,height=5.2cm,keepaspectratio]{horizon_step}{3}
    \end{column}
    \begin{column}{0.30\textwidth}
      \small$\Pr(\text{success}) = (1-\epsilon)^{H}$\\[0.1cm]
      {\footnotesize $\epsilon$: error per step; $H$: turns}\\[0.3cm]
      \wstep{2}{$\epsilon = 1 - \val{target}^{1/\val{h_long}} = \val{eps_req_pct}\%$}
      \wstep{2}{\val{eps_ratio}$\times$ lower than \val{step_err_pct}\%}
      \wstep{3}{$\epsilon H \approx \val{eps_h}$ holds \val{target_pct}\% at any $H$}
    \end{column}
  \end{columns}
  \mlsyscredit{Introduction, equation 1 and Napkin Math 1.1 (independent, unchecked steps; illustrative \val{step_ok_pct}\%). Values: calcs.py.}
  \note{[2.5 min] [CORE] Worked example. Step 1: with independent steps, no check, and no recovery, all $H$ must succeed: $\val{step_ok}^{\val{h_short}} = \val{p_short_pct}$ percent, $\val{step_ok}^{\val{h_long}} = \val{p_long_pct}$ percent. Step 2: to hold \val{target_pct} percent at \val{h_long} turns, $\epsilon = 1 - \val{target}^{1/\val{h_long}} = \val{eps_req}$, \val{eps_req_pct} percent per step, \val{eps_ratio} times lower. For small $\epsilon$, $(1-\epsilon)^H \approx e^{-\epsilon H}$, so the product $\epsilon H \approx \val{eps_h}$ sets reliability at any horizon. Step 3: doubling the unchecked horizon to \val{h_ext} squares success, \val{p_ext_pct} percent. The cheaper knob is the unchecked horizon: a check that catches an error and restores state cuts the run into short segments; with perfect checks every \val{h_short} turns, each segment succeeds \val{p_short_pct} percent of the time and is retried. ENGAGE: in the astropy run, which check shortened the unchecked horizon? (Expected: the interpreter caught the bad edit one turn later.) WARN: real errors correlate and models sometimes recover; the shape of the curve can change while the lesson holds. FLEX: [CORE].}
\end{frame}

\begin{frame}{Halve the token price, or raise acceptance?}
  \mlsysactivity{Estimate}{3 min}
  \begin{minipage}[c]{0.40\textwidth}
      \begin{mlsyscard}{mlsysaccent}
        \footnotesize\textbf{Given: one repair attempt}\\[0.03cm]
        Tokens \$\val{c_tokens}; sandbox, tests, review \$\val{c_other}; accepted \val{p_acc_pct}\%; retry until accepted.\\[0.03cm]
        1: halve the token price. 2: acceptance to \val{p_better_pct}\%.
      \end{mlsyscard}
  \end{minipage}\hfill
  \begin{minipage}[c]{0.56\textwidth}
      \centering\mlsyssteps[width=\textwidth,height=2.9cm,keepaspectratio]{costbars_step}{2}
  \end{minipage}\\[0.2cm]
  \small
  \wstep{2}{$C_{\text{eff}} = \mathbb{E}[C_{\text{attempt}}]/p_{\text{acc}} = \val{c_attempt}/\val{p_acc} = \$\val{c_eff}$}
  \wstep{2}{Half the token price: $\val{c_cheap}/\val{p_acc} = \$\val{c_eff_cheap}$.\quad Acceptance \val{p_better_pct}\%: $\val{c_attempt}/\val{p_better} = \$\val{c_eff_better}$.}
  \wstep{3}{Break-even acceptance for the price cut: $\val{p_acc}/(1 - \val{token_cut} \times \val{token_share}) = \val{p_breakeven}$}
  \mlsyscredit{Introduction, equation 2 and Napkin Math 1.2 (illustrative prices). Values: calcs.py.}
  \note{[3 min] [CORE] ENGAGE: two minutes in pairs; each pair picks an option before the reveal. Step 2: with retries until acceptance, attempts are geometric with mean $1/p_{\text{acc}} = \val{attempts}$, so cost per accepted task is $C_{\text{eff}} = \mathbb{E}[C_{\text{attempt}}]/p_{\text{acc}}$: \$\val{c_attempt} / \val{p_acc} = \$\val{c_eff}. Halving the token price makes the attempt \$\val{c_cheap} and gives \$\val{c_eff_cheap}; raising acceptance to \val{p_better_pct} percent gives \$\val{c_eff_better}. Step 3: tokens are \val{token_share_pct} percent of the attempt, so the price cut is worth the same as acceptance rising to \val{p_breakeven}. The levers act on different amounts: a price cut removes part of one component, while acceptance divides the whole attempt, sandbox and review included. WARN: students who chose the price cut reason from the provider's bill; the right bill charges every attempt. A mechanism pays only if $C'/p' < C/p$. FLEX: [CORE].}
\end{frame}

\begin{frame}{One real repair agent paid \$\val{run_per_acc} per accepted task}
  \centering
  \safeimg[width=0.90\textwidth,height=3.6cm,keepaspectratio]{realcost.pdf}\\[0.15cm]
  {\small $\$\val{run_mean} / \val{run_acc_frac} = \$\val{run_per_acc}$. Charging only the accepted runs understates the cost \val{run_understate}$\times$.}
  \mlsyscredit{SWE-bench Verified run as before: \$\val{run_total} over \val{run_tasks} attempts, \val{run_acc} accepted; per accepted task and the ratio derived in calcs.py.}
  \note{[2 min] [OPTIONAL] NARRATE: equation 2 on real numbers. The run spent \$\val{run_total} on \val{run_tasks} single attempts, a mean of \$\val{run_mean}. The hidden tests accepted \val{run_acc}, \val{run_acc_pct} percent, so each accepted repair cost \$\val{run_mean} / \val{run_acc_frac} = \$\val{run_per_acc3}. The accepted runs alone averaged \$\val{run_winners}; reporting that as the cost per repair hides the \val{run_rej} failures and understates by \val{run_understate2} times. The chapter's illustrative case understates by \val{c_understate} times. ENGAGE: what is missing from this \$\val{run_per_acc}? (Expected: sandbox compute, the test runs, and any human review; the bill counts only model tokens.) WARN: with one attempt per task, the per-attempt form and the cohort form give the same number here. Lecture 19 owns the cohort form, which charges every initiated task, including canceled ones. FLEX: cut if short.}
\end{frame}

\begin{frame}{A model-directed loop pays only where it raises acceptance}
  \centering
  \safeimg[width=0.92\textwidth,height=3.5cm,keepaspectratio]{workflow.pdf}\\[0.15cm]
  {\small The loop earns its place when $\dfrac{C_{\text{loop}}}{p_{\text{loop}}} < \dfrac{C_{\text{workflow}}}{p_{\text{workflow}}}$ under the same contract.}
  \mlsyscredit{Introduction, \S1.6.1.}
  \note{[1.5 min] [OPTIONAL] NARRATE: in the repair, cloning, installing dependencies, and starting the linter never change order: a fixed workflow does them with no tokens and no new failure points. Diagnosing an unknown race does change with what the last step found, so a model-directed loop belongs there. A fixed workflow can still branch, retry, and run in parallel; its control graph is fixed in advance. ENGAGE: a known flaky test fails; loop or retry rule? (Expected: a retry rule.) WARN: compare both designs under the same contract and the same acceptance checks, or the comparison measures the grader. Approval policy is separate from control structure: a workflow can require approval, and a loop can run unprompted in a sandbox. Land the segment aloud: horizon multiplies risk, acceptance divides cost, and both equations price every mechanism. FLEX: cut if short; say the land line on the previous frame.}
\end{frame}

% =============================================================================
% SECTION: ONE TURN
% =============================================================================
\mlsysnote{[0.5 min] [CORE] Fifth: what happens inside one turn, and which parts belong to the model?}
\mlsysdivider{What happens inside one turn?}{Introduction, \S1.7}
\section{One Turn}

\begin{frame}{Every turn passes six phases, and the model owns one}
  \centering
  \mlsyssteps[width=0.95\textwidth,height=5.9cm,keepaspectratio]{sixphase_step}{3}
  \mlsyscredit{Redrawn from Introduction, figure 5.}
  \note{[2.5 min] [CORE] Step 1: the runtime assembles context from the contract, the workspace, and recent observations, within token and dollar budgets. The model returns a proposal; the runtime treats a cut-off or malformed reply as a failed call and never executes it. The runtime validates the proposal against the tool's schema and the task's grant. Step 2: only an authorized request reaches a tool, which runs in isolation with timeouts. The runtime records the outcome as evidence, then assesses the task against the contract and starts the next turn. A call whose reply is lost has an unknown outcome; Lecture 7 owns settling it before any retry. Step 3: a rejected proposal returns as error feedback without touching a tool. Waiting keeps the task alive while the model is idle: holding its serving memory through a long test is the tool-wait problem, priced in Lecture 5. Every turn ends accepted, failed, unresolved, or escalated; an exhausted budget counts as unresolved, because the evidence neither accepts nor rejects the work. ENGAGE: which phase would have stopped the hook's delete? (Expected: 3.) FLEX: [CORE].}
\end{frame}

\begin{frame}{The real repair, phase by phase}
  \footnotesize
  \begin{tabular}{@{}P{0.9cm}P{9.6cm}P{3.6cm}@{}}\toprule
    \textbf{Turn} & \textbf{Recorded proposal or observation} & \textbf{Phase}\\\midrule
    11 & {\ttfamily \textcolor{RedLine}{$\times$} mask * constant: unsupported ...} & 5: bug evidence\\
    13 & {\ttfamily sed -i '521a\textbackslash{} elif operand ...'} & 2 to 4: patch\\
    14 & {\ttfamily \textcolor{RedLine}{IndentationError: expected an ...}} & 5: caught next turn\\
    16 & {\ttfamily git checkout .../ndarithmetic.py} & 4: restore\\
    19 & {\ttfamily \textcolor{GreenLine}{mask * no mask: OK}} & 5: visible check\\
    28 & {\ttfamily echo MICRO\_SWE\_AGENT\_FINAL\_OUTPUT} & 6: model says done\\
    after & {\ttfamily test\_nddata\_bitmask\_arithmetic} + \val{hook_p2p} & accepted outside\\\bottomrule
  \end{tabular}
  \mlsyscredit{Verbatim excerpts (shortened where marked), astropy\_\_astropy-14995 trajectory, SWE-bench experiments; resolution from per\_instance\_details.json.}
  \note{[2 min] [CORE] NARRATE: the chapter walks an illustrative race-condition repair; this is the same anatomy in a recorded run. Turn 11: the reproduction script shows the crash, evidence in phase 5. Turn 13: the model proposes a one-line insertion with \code{sed}; the harness runs it. Turn 14: the interpreter rejects the edited file, the same role the compiler's unclosed-bracket error plays in the chapter: a check that catches the error one turn after it happened. Turn 16: the model restores the file and patches again at turn 17. Turn 19: the agent's own script passes, level-3 evidence. Turn 28: the model declares done. Acceptance came after the loop, from sealed tests. ENGAGE: which phase in this harness is thinnest? (Expected: phase 3; it checks only that the reply holds one command.) FLEX: [CORE].}
\end{frame}

\begin{frame}[fragile]{A bare loop stops on the model's word; a harness stops on limits}
  \begin{columns}[T]
    \begin{column}{0.52\textwidth}
      {\small\textbf{Naive loop (chapter, Listing 1)}}
\begin{semiverbatim}\fontsize{8.5}{10}\selectfont\bfseries
prompt = task_description
while True:
    response = generate(prompt)
    command = parse_command(response)
    output = \textcolor{RedLine}{run_with_host_permissions}(
                 command)
    prompt += format_observation(...)
    if \textcolor{RedLine}{"TASK_COMPLETE" in response}:
        break
\end{semiverbatim}
    \end{column}
    \begin{column}{0.46\textwidth}
      {\small\textbf{A real harness's budget (config)}}
\begin{semiverbatim}\fontsize{8.5}{10}\selectfont\bfseries
step_limit: int = 0
cost_limit: float = 3.0
wall_time_limit_seconds: int = 0
max_consecutive_format_errors: int = 3
\end{semiverbatim}
      {\footnotesize Exit statuses: \code{Submitted}, \code{LimitsExceeded},\\ \code{TimeExceeded}, \code{RepeatedFormatError}}
    \end{column}
  \end{columns}
  \vspace{0.2cm}
  {\small Outcomes: accepted, failed, \textbf{unresolved} (budget spent), escalated. A silent stop is a defect.}
  \mlsyscredit{Left: Introduction, Listing 1. Right: mini-swe-agent 2.4.6, the current version (the 2025 runs used 1.0.0), agents/default.py (commit 04d809c, MIT).}
  \note{[2 min] [OPTIONAL] NARRATE: the chapter's naive loop runs every command with host permissions, appends unbounded output to the prompt, and stops when the model prints a word. If the model cycles on a broken patch it never stops; if the process crashes after writing a file but before recording the test, a restart runs against a dirty tree. The real harness adds budgets: steps, dollars (\$3 by default), wall time, repeated format errors, each with its own exit status. ENGAGE: when the dollar limit trips, is the task failed? (Expected: unresolved: the evidence neither accepts nor rejects it; save the diff and escalate.) WARN: limits stop the bleeding but do not settle effects or survive crashes. LINK: Lecture 10 makes the harness own the decision to continue; Lecture 11 records intent before each effect so a crash cannot repeat it. Land the segment aloud. FLEX: cut if short.}
\end{frame}

% =============================================================================
% CLOSE
% =============================================================================
\begin{frame}{Back to the hook: why did this repair count?}
  \mlsysactivity{Retrieval}{2.5 min}
  \begin{columns}[c]
    \begin{column}{0.38\textwidth}
      \centering\safeimg[width=\textwidth,height=2.4cm,keepaspectratio]{issue_14978.png}\\[0.1cm]
      {\normalsize\only<1|handout:0>{Name the loop, the exposure, the evidence, and the price.}}
    \end{column}
    \begin{column}{0.60\textwidth}
      \small
      \uncover<2->{\begin{tabular}{@{}P{1.8cm}P{6.0cm}@{}}\toprule
        \textbf{Loop} & a runtime ran \val{hook_calls} proposals; none had authority\\[0.04cm]
        \textbf{Exposure} & $H = \val{hook_calls}$, $S_2$, $A_1$: a container copy\\[0.04cm]
        \textbf{Evidence} & a sealed test plus \val{hook_p2p} old tests\\[0.04cm]
        \textbf{Price} & \$\val{hook_cost} here; \$\val{run_per_acc} per accepted task over \val{run_tasks}\\\bottomrule
      \end{tabular}}\\[0.2cm]
      \uncover<3->{\small The same delete is harmless in this copy and fatal on a developer's only workspace.}
    \end{column}
  \end{columns}
  \mlsyscredit{As on the hook; values from calcs.py.}
  \note{[2.5 min] [CORE] ENGAGE: 60 seconds of silent writing first. Step 2: the four answers. The model proposed \val{hook_calls} commands and held no authority: the loop ran each in a container and recorded what came back. The task sat at \val{hook_calls} turns, a working copy, and workspace authority, so letting the model act freely risked only a discarded copy. Acceptance came from a test the agent never saw; its sentence counted for nothing, and \val{run_rej} equally confident runs failed that test. The price: \$\val{hook_cost} for this task, \$\val{run_per_acc} per accepted task once the failures are charged. Step 3: the hook's second question. The model could not tell the two workspaces apart; only the runtime's environment could. That is the answer to the driving question: a correct action fails a task when nothing closes the loop around it. FLEX: [CORE].}
\end{frame}

\begin{frame}{Takeaways}
  \thumbline{thumb_loop_step2.pdf}{Agency belongs to the closed loop; the model proposes, and the runtime decides and records.}\\[0.1cm]
  \thumbline{thumb_hsa_step3.pdf}{Horizon, state, and authority locate a task; authority is the consequence in its environment.}\\[0.1cm]
  \thumbline{thumb_ladder_step2.pdf}{Bounds close below the model, evidence above it, both written into a contract first.}\\[0.1cm]
  \thumbline{thumb_horizon_step3.pdf}{$(1-\epsilon)^H$ prices unchecked turns; $\mathbb{E}[C_{\text{attempt}}]/p_{\text{acc}}$ prices every mechanism.}\\[0.1cm]
  \thumbline{thumb_sixphase_step3.pdf}{Each turn passes six phases; only model invocation belongs to the model.}
  \note{[1 min] [CORE] One line per segment. Ask students which slide taught each line. If short on time, read the first and the third.}
\end{frame}

\begin{frame}{Open questions}
  \normalsize
  \begin{itemize}\setlength{\itemsep}{0.3cm}
    \item What will limit delegated tasks first: model accuracy, cost per accepted task, or the evidence needed to accept the work?
    \item When an agent's action cannot be undone, who should approve it, and how long may approval take?
    \item If delegated tasks keep doubling every few months, what must change in the runtime rather than in the model?
  \end{itemize}
  \note{[1 min] [OPTIONAL] No single answer. Pick one for discussion or assign one as a short written response. The first ties the METR trend to the two equations; the second is the remediation agent's problem; the third asks which checks must shorten the unchecked horizon as tasks lengthen.}
\end{frame}

\begin{frame}{Read and explore}
  \begin{columns}[c]
    \begin{column}{0.74\textwidth}
      \small
      \begin{itemize}\setlength{\itemsep}{0.02cm}
        \item \textbf{Read:} Introduction, \S1.1 to \S1.9; Checkpoint 1.1; Napkin Math 1.1 and 1.2 by hand.
        \item Saltzer and Schroeder, ``Protection of Information,'' 1975.
        \item Saltzer, Reed, and Clark, ``End-to-End Arguments,'' 1984.
        \item Avizienis et al., ``Dependable and Secure Computing,'' 2004.
        \item Yao et al., ReAct, 2023; Jimenez et al., SWE-bench, 2024.
        \item Kwa et al., ``Long Software Tasks,'' 2025.
      \end{itemize}
    \end{column}
    \begin{column}{0.24\textwidth}
      \centering\safeimg[width=\textwidth,height=3.0cm,keepaspectratio]{norbert_wiener.jpg}\\[0.05cm]
      {\footnotesize Norbert Wiener, \textit{Cybernetics}, 1948: feedback from the environment.}
    \end{column}
  \end{columns}
  \mlsyscredit{Photo: Konrad Jacobs, CC BY-SA 2.0 de, Wikimedia Commons.}
  \note{[1 min] [CORE] The chapter first, with Checkpoint 1.1 and the two notebooks done by hand. Saltzer and Schroeder give least privilege and complete mediation; the end-to-end argument explains why only an endpoint that knows the task can check it; Avizienis et al. name the unsignaled failure. ReAct is the reason-act-observe loop; SWE-bench is the benchmark behind today's real run; Kwa et al. measured the horizon trend. The loop itself is old: Wiener's 1948 \textit{Cybernetics} made feedback from the environment the basis of control, and he warned in 1960 that a machine acting too fast to stop must carry exactly the purpose we intend. FLEX: students who want to read a trajectory can open any file in the SWE-bench experiments repository.}
\end{frame}

\begin{frame}[c]
  \begin{columns}[c]
    \begin{column}{0.34\textwidth}
      \centering\mlsyskitfig[width=\textwidth,height=4.2cm,keepaspectratio]{stack_model}
    \end{column}
    \begin{column}{0.62\textwidth}
      {\small\textcolor{lightgray}{NEXT LECTURE \textbar{} THE FOUNDATION MODEL}}\\[0.5cm]
      {\LARGE\bfseries What does one model call return, and what does it cost?}
    \end{column}
  \end{columns}
  \note{[0.5 min] [CORE] Of the six phases, only the second belongs to the model, and the other five depend on knowing exactly what it delivers: a complete tool request, one cut off midway, or nothing. Lecture 2 follows one inspection call of the repair agent from request to bill.}
\end{frame}

\end{document}
